% !PRINTED 201
\weberpage{201}

\begin{center}
§. 65.

\gesperrt{Auflösung der biquadratischen Gleichung.}
\end{center}

Wir wollen eine Normalform durch eine lineare Substitution
mit der Determinante 1 herstellen, zu deren Ableitung wir die
gegebene biquadratische Form in lineare Factoren zerlegt annehmen
:
\weq{1}{f(x,\,y) = a_0\,(x - \alpha y)\,(x - \beta y)\,(x - \gamma y)\,(x - \delta y).}
Wir setzen
\weq{2}{\begin{gathered}
x\,\sqrt{\alpha - \beta} = \beta\,\xi - \alpha\,\eta, \quad r = 1,\\
y\,\sqrt{\alpha - \beta} = \xi - \eta,
\end{gathered}}
also
\weq{3}{\begin{gathered}
x - \alpha y = -\sqrt{\alpha - \beta}\,\xi\\
x - \beta y = -\sqrt{\alpha - \beta}\,\eta\\
(x - \gamma y)\,\sqrt{\alpha - \beta} = (\beta - \gamma)\,\xi - (\alpha - \gamma)\,\eta\\
(x - \delta y)\,\sqrt{\alpha - \beta} = (\beta - \delta)\,\xi - (\alpha - \delta)\,\eta,
\end{gathered}}
und dadurch geht $f(x,\,y)$ über in
\weq{4}{F(\xi,\,\eta) = a_1'\,\xi^3\eta + a_2'\,\xi^2\eta^2 + a_3'\,\xi\eta^3.}

Die Coëfficienten $a_1', a_2', a_3'$ sind nach (3) leicht durch die
$\alpha, \beta, \gamma, \delta$ auszudrücken:
\weq{5}{\begin{gathered}
a_1' = \phantom{-{}} a_0\,(\beta - \gamma)\,(\beta - \delta)\\
a_3' = \phantom{-{}} a_0\,(\alpha - \gamma)\,(\alpha - \delta)\\
a_2' = -a_0\,[(\alpha - \gamma)\,(\beta - \delta) + (\alpha - \delta)\,(\beta - \gamma)]
\end{gathered}}

Durch Vertauschung der Wurzeln können wir die Normalform
 (4) auf sechs verschiedene Arten herstellen, die aber nur
drei verschiedene Werthe von $a_2'$ liefern. Diese drei Werthe
sind, wenn wie oben
\weq{6}{\begin{gathered}
U = (\alpha - \beta)\,(\gamma - \delta),\\
V = (\alpha - \gamma)\,(\delta - \beta),\\
W = (\alpha - \delta)\,(\beta - \gamma),\\
U + V + W = 0
\end{gathered}}
gesetzt wird,
%CONT>>
